\section{Scalar facts for the critical tanh variance chain}
\label{sec:echo-scalar-preliminaries}
Write
\[
 F(q)=\E\tanh^2(\sqrt q Z),\qquad Z\sim N(0,1),\qquad
 \chi(q)=\E\tanh'(\sqrt q Z)=1-F(q).
\]
The following estimates also make the constants in the covariance schedules explicit in their dependence on a fixed input dimension.

\begin{lemma}[Variance drift and derivative products]
\label{lem:echo-scalar-facts}
The function $F$ is continuous and strictly increasing on $(0,\infty)$, with $0\le F'(q)\le1$. For absolute $c,c'>0$ and $0<q\le1$,
\[
 c'q\le F(q),\qquad q-2q^2\le F(q)\le q-cq^2.
\]
Fix $q_0\in(0,1)$ and put $q_{j+1}=F(q_j)$. There are positive constants depending only on $q_0$ such that
\[
 \frac c{j+1}\le q_j\le\frac C{j+1},\qquad
 c\sqrt{q_m/q_0}\le\prod_{j=0}^{m-1}(1-q_{j+1})
                \le C\sqrt{q_m/q_0}.
\]
In particular the product is comparable to $(m+1)^{-1/2}$.
\end{lemma}
\begin{proof}
Strict monotonicity and continuity follow pointwise from $\tanh^2(\sqrt q z)$ and bounded convergence. The heat-semigroup differentiation identity, justified by bounded derivatives, gives
\[
 F'(q)=\tfrac12\E(\tanh^2)''(\sqrt q Z).
\]
With $u=\tanh^2 t$, the half second derivative is $1-4u+3u^2\le1$ for $0\le u\le1$. Monotonicity supplies the lower derivative bound.

For $0\le t\le1$, concavity gives $\tanh t\ge t/2$, and hence
\[
 t-\tanh t=\int_0^t\tanh^2 s\,ds\ge t^3/12.
\]
Thus $t^2-\tanh^2t\ge t^4/12$. Restricting the Gaussian integral to $|Z|\le1$ proves the upper drift bound and the lower linear bound. For the other drift inequality use
$\tanh^2t\ge t^2-(2/3)t^4$ for all real $t$. If its right side is positive, $|t|\le\sqrt{3/2}$ and the inequality follows by squaring $\tanh t\ge t-t^3/3$ for $t\ge0$; this last bound follows by integrating $\tanh'(t)=1-\tanh^2t\ge1-t^2$. Where the right side is nonpositive the claim is immediate. Gaussian averaging gives $F(q)\ge q-2q^2$.

The reciprocal increment is bounded above and below by positive constants:
\[
 c\le\frac1{q_{j+1}}-\frac1{q_j}
 =\frac{q_j-F(q_j)}{q_jF(q_j)}\le2/c'.
\]
Summation proves the bounds on $q_j$.

Taylor expansion with a bounded remainder gives
\[
 F(q)=q-2q^2+\tfrac{17}{3}q^3+O(q^4).
\]
For example expand $\tanh^2t=t^2-(2/3)t^4+(17/45)t^6+O(t^8)$ with a globally valid bounded-remainder constant and average the Gaussian moments. Consequently
\[
 J(q):=\frac{q[1-F(q)]^2}{F(q)}
       =1-\tfrac23q^2+O(q^3).
\]
Every $J(q_j)$ is positive. Since $\sum_jq_j^2<\infty$, the product $\prod_{j<m}J(q_j)$ is bounded above and below by positive constants independent of $m$: its tail logarithm is absolutely summable and its finite initial segment is positive. The exact telescoping identity
\[
 \prod_{j=0}^{m-1}(1-q_{j+1})^2
   =\frac{q_m}{q_0}\prod_{j=0}^{m-1}J(q_j)
\]
proves the derivative-product estimate.
\end{proof}
